%! Author = sriadityavedantam
%! Date = 4/7/26

\section{Jouanolou's Trick}

\begin{theorem}[Jouanolou's Trick]{
    There exists an affine variety $X$ together with a regular map
    \[
        X \to \pp^n
    \]
    such that all fibers are affine varieties.
}
    \begin{definition}[Fiber]
        The fiber over a point $p$ is the preimage $f^{-1}(p)$.
    \end{definition}
    We construct the example for $\pp^1$.
    Recall that $\pp^1$ parametrizes lines through the origin in $\aa^2$.
    Fix the line $l \coloneqq \{y = 0\}\subset \aa^2$.
    Consider linear maps $M:\aa^2\to\aa^2$ with image equal to a line.
    Such maps are represented by $2\times 2$ matrices of rank $1$.
    For example,
    \[
        M =
        \begin{bmatrix}
            1 & 0\\
            0 & 0
        \end{bmatrix}
        \quad\text{gives}\quad
        M\begin{bmatrix}x\\y\end{bmatrix}
        =
        \begin{bmatrix}x\\0\end{bmatrix}.
    \]
    \begin{definition}[Rank]
        The rank is the dimension of the image, equivalently the number of linearly independent columns.
    \end{definition}
    Assume $M^2 = M$.
    Then $M$ is a projection, and we obtain a decomposition
    \[
        \aa^2 \cong \ker(M)\oplus \ima(M).
    \]
    The eigenvalues of $M$ satisfy
    \[
        x^2 - x = 0,
    \]
    so they are $0$ and $1$.
    The $0$-eigenspace is $\ker(M)$ and the $1$-eigenspace is $\ima(M)$.
    Now define $X$ to be the set of $2\times 2$ matrices
    \[
        M =
        \begin{bmatrix}
            a & b\\
            c & d
        \end{bmatrix}
    \]
    such that
    \[
        M^2 = M \quad\text{and}\quad \rank(M) = 1.
    \]
    We view $X$ as a subset of $\aa^4$ with coordinates $(a,b,c,d)$.
    The condition $M^2 = M$ gives the polynomial equations
    \begin{align*}
        a^2 + bc &= a,\\
        ab + bd &= b,\\
        ac + dc &= c,\\
        bc + d^2 &= d.
    \end{align*}
    To ensure $\rank(M)=1$, we impose
    \[
        \det(M) = ad - bc = 0,
    \]
    and exclude the zero matrix $M=0$.
    Thus $X$ is a quasi-affine variety (an open subset of an affine variety).
    Now define a map
    \[
        \pi : X \to \pp^1
    \]
    by
    \[
        M \longmapsto \ima(M).
    \]
    Since $\ima(M)$ is a $1$-dimensional subspace of $\aa^2$, it determines a point in $\pp^1$.
    This map is regular.
    For a fixed line $L\subset \aa^2$, the fiber
    \[
        \pi^{-1}(L)
    \]
    consists of all rank $1$ idempotent matrices with image $L$.
    Such matrices are determined by a choice of complementary kernel, and this parameter space is affine.
    Hence all fibers are affine varieties.
    This constructs an affine variety $X$ mapping to $\pp^1$ with affine fibers, which is the desired statement.
\end{theorem}